\documentclass[10pt]{article}
\usepackage{resume_shared}

\begin{document}
\fontsize{9.8pt}{11.7pt}\selectfont

\ResumeName{Sri Ravi Kumar Birada}
\ResumeContact{Hyderabad, India}{(+91) 7200999942}{ravikumar31birada@gmail.com}{www.linkedin.com/in/ravi-20043}

\section{Summary}
AI consultant with 7+ years advising enterprises on AI-led workflow transformation, SaaS product modernization, governance-aware adoption, and business-case development. Strong at facilitating discovery, applying value stream mapping, shaping use cases from ideation to MVP and evaluation, and translating technical capabilities into executive-ready narratives. Experience spans AI roadmap definition, 0 to 1 capability shaping, agentic AI for RFP automation, and enterprise pilot readiness.

\section{Professional Experience}
\RoleEntry{AI Product Manager - Enterprise AI Products}{HCLTech}{2023 -- Present}
\begin{itemize}
  \item Formulated an AI-specific transformation roadmap across 3 enterprise SaaS product modules by mapping 20+ workflow activities and prioritizing 5 use cases over a 12-month horizon.
  \item Introduced an agentic AI RFP response workflow that ingests customer RFPs and enterprise capability / technology-stack data, adds validation gates, and improves readiness of final submission documents.
  \item Shaped a 0 to 1 enterprise SaaS capability by synthesizing inputs from 6+ customer accounts and positioning 8 MVP feature areas across platform-extension and standalone adoption paths.
  \item Mobilized end-to-end onboarding of 3 GenAI partners by coordinating stakeholders, identifying 3 POC / pilot accounts, and defining 5 pilot success measures for enterprise adoption.
  \item Diagnosed workflow-modernization opportunities across field operations, commercial operations, exception handling, knowledge access, and document generation, resulting in clearer transformation pathways.
  \item Balanced AI ambition with enterprise controls by incorporating governance needs, integration complexity, adoption risks, and operating-model constraints into client-facing recommendations.
\end{itemize}

\RoleEntry{Product Manager - Enterprise SaaS Products}{HCLTech}{May 2022 -- 2023}
\begin{itemize}
  \item Captured multi-account customer inputs and converted them into modular SaaS requirements across customer operations, workflow orchestration, reporting, and commercial governance.
  \item Commercialized pay-as-you-grow platform narratives by connecting product capability, buyer pain points, implementation sequencing, and measurable time-to-value expectations.
  \item Advanced operational delivery for an enterprise product deployment supporting 300K+ SKUs and 52+ distributed business entities, contributing to launch readiness within 21 days.
\end{itemize}

\MiniHeading{Selected Consulting Outcomes}
\begin{itemize}
  \item AI transformation roadmap: converted 20+ workflow activities into 5 AI use cases, 3 enterprise SaaS product-module priorities, and a 12-month execution view.
  \item Agentic RFP acceleration: framed an AI workflow for customer RFP ingestion, enterprise knowledge retrieval, validation, and final-document sanity review.
  \item 0 to 1 SaaS capability: translated 6+ customer-account inputs into 8 MVP feature areas across platform extension and standalone adoption paths.
  \item Partner pilot advisory: connected 25+ partner evaluations with 3 onboarded partners, POC-account selection, 5 success metrics, and enterprise adoption readiness.
\end{itemize}

\RoleEntry{Network Engineer}{Ericsson Global India Limited}{Oct 2016 -- Jan 2019}
\begin{itemize}
  \item Reduced deliverable turnaround time by 15\% by automating report validation workflows with Python.
  \item Cut server infrastructure requirements by 40\% through performance improvements in analysis tooling and workflow design.
  \item Received the Ericsson Power Award in Q4 2017 for automation impact in Single Site Verification reporting.
\end{itemize}

\section{Education}
\InlineSectionText{MBA, MDI Gurgaon (2022)}
\InlineSectionText{B.Tech, Electronics and Communication Engineering, Amrita School of Engineering (2016)}

\section{Certifications}
\InlineSectionText{OpenAI AI Technical Practitioner (2026) \textbar\ Google Certified GenAI Leader (2025) \textbar\ GenAI Prompt Engineer L2 (2025) \textbar\ Certified Scrum Product Owner (CSPO) (2024)}

\section{Skills}
\SkillGroup{Consulting and Strategy}{AI Consulting, Enterprise AI Strategy, Transformation Advisory, Use-Case Discovery, Workshop Facilitation, Stakeholder Alignment, Business Case Development, 0 to 1 Module Shaping}
\SkillGroup{Enterprise AI}{AI Governance, Responsible AI, Workflow Transformation, GenAI Strategy, Knowledge Search, Agentic Workflow Design}
\SkillGroup{Tools and Methods}{Value Stream Mapping, Partner Onboarding, Pilot Success Definition, Power BI, Python, Jira, Cross-Functional Solutioning}

\end{document}
